\section{What a Subgraph Owes the Router}
\label{sec:ch07-what-a-subgraph-owes}
\index{subgraph!contract|(}%

A subgraph is an ordinary GraphQL service that has agreed to two extra fields.
The Apollo Federation subgraph specification names them, along with the types
they need~\autocite{apollo2026subgraphspec}:

\begin{graphqlsdl}
extend type Query {
  _service: _Service!
  _entities(representations: [_Any!]!): [_Entity]!
}

type _Service { sdl: String! }
union _Entity
scalar _Any
\end{graphqlsdl}

\index{HotChocolate!AddApolloFederation@\code{AddApolloFederation()}}%
In HotChocolate 16.6.0 you get all of it from one call. The Catalog subgraph's
entire federation configuration is this:

\begin{minted}{csharp}
builder.AddGraphQL()
    .AddApolloFederation()
    .AddWireCatalog();
\end{minted}

\code{AddWireCatalog()} is the source generator's registration method, the same
mechanism chapter~\ref{ch:hotchocolate-quickly} introduced.
\code{AddApolloFederation()} is the new line, and here is what it puts into the
schema:

\begin{graphqlsdl}
type Query {
  products: [Product!]!
  productById(id: ID!): Product
  _service: _Service!
  _entities(representations: [_Any!]!): [_Entity]!
}

type _Service { sdl: String! }
union _Entity = Product
scalar FieldSet
scalar _Any

directive @key(fields: FieldSet!, resolvable: Boolean = true) repeatable on
  | OBJECT
  | INTERFACE
\end{graphqlsdl}

\index{federation!Entity union@\code{\_Entity} union}%
\code{\_Entity} is a union of exactly the types carrying \code{@key}, which here
is one type. That is worth holding on to: the union is generated from the keys,
so a type without a key can never appear in it, and section~\ref{sec:ch07-representations}
is what happens when the router is told otherwise.

\subsection*{The version HotChocolate speaks is not the version Apollo publishes}

\index{federation!specification version}%
The first line of every subgraph schema declares which federation specification
it is written against:

\begin{graphqlsdl}
schema
  @link(
    url: "https://specs.apollo.dev/federation/v2.6"
    import: ["@key", "@tag", "FieldSet"]
  ) {
  query: Query
}
\end{graphqlsdl}

Version 2.6, against a specification that is at v2.15. That is not a typo and it
is not a setting I forgot. \code{AddApolloFederation()} takes an optional
\code{FederationVersion}, and the enum in HotChocolate 16.6.0 stops at 2.7:

\begin{minted}{csharp}
Federation26 = 2_6,
Federation27 = 2_7,
// default to latest-1
Default = Federation26,
Latest = Federation27
\end{minted}

\index{HotChocolate!FederationVersion@\code{FederationVersion}}%
The comment is theirs. The default is deliberately one behind the newest version
the library supports, which is a reasonable posture for a library and a
surprising one to meet in your own schema output. Pass
\code{FederationVersion.Latest} and you get 2.7; there is no way to ask for
anything above it.

In practice this has cost me nothing. Every directive
chapter~\ref{ch:the-federation-model} covers was in the specification by v2.3,
the composer accepted a v2.6 subgraph without a word, and the router planned
against it correctly. But it is a ceiling, and it is worth knowing it is there
before you go looking for something the specification added after v2.7.

\subsection*{\code{\_service}, and who actually reads it}

\index{federation!service field@\code{\_service} field}%
\code{\_service} returns the subgraph's own schema as a string. Ask Catalog for
it and you get back the SDL above, directives and all. The specification is
particular about this: the string \enquote{must include all uses of all
federation-specific directives, such as
\code{@key}}~\autocite{apollo2026subgraphspec}. A schema printed without its
\code{@key} directives would compose into a graph with no entities in it.

The natural assumption is that the router calls this at startup. It does not.
Start both subgraphs, start the router, wait for it to report itself ready, then
count the occurrences of \code{\_service} in each subgraph's console. Zero on
both. The router never asked.

That follows from the third startup warning of
section~\ref{sec:ch07-three-processes}. A statically configured router got its
schema from a file, so there is nothing left to ask about. \code{\_service} is
the \emph{composer's} field, not the router's.

\index{composition!introspection mode}%
And the composer really does use it. An \code{introspection} block in
\code{graph.yaml}, where the schema file would otherwise go, points it at a live
endpoint instead. Run the compose command again and this arrives:

\begin{minted}[fontsize=\footnotesize]{text}
User-Agent: node
RequestBody: {"query":"\n  {\n    _service{\n      sdl\n    }\n  }\n","variables":{}}
\end{minted}

The whitespace is the composer's own document, sent as it was written. Nothing
clever, nothing proprietary: the tool asks the same question you can ask from
Postman, and the two committed schema files in this sample were produced by
asking it.

\subsection*{HotChocolate's own description of \code{\_service} is wrong}

\index{HotChocolate!Service type description@\code{\_Service} type description}%
Every schema HotChocolate generates carries this description on \code{\_Service},
and I am quoting it in full because the interesting part is at the end:

\begin{minted}[fontsize=\footnotesize,breaklines]{text}
This type provides a field named sdl: String! which exposes the SDL of the
service's schema. This SDL (schema definition language) is a printed version of
the service's schema including the annotations of federation directives. This
SDL does not include the additions of the federation spec.
\end{minted}

The SDL it returns includes \code{\_service}, \code{\_entities}, \code{\_Any},
\code{\_Entity}, \code{FieldSet}, \code{@key} and \code{@link}. Every one of
those is an addition of the federation spec.

The behaviour is correct and the sentence is not. Apollo's specification
forbids including those definitions only when a subgraph supports Federation 1;
for Federation 2 it says outright that \code{sdl} may include
them~\autocite{apollo2026subgraphspec}. So the description is a true statement
about a rule that does not apply here, shipped as documentation of something
that does. Read the output, not the docstring. That is the whole of
chapter~\ref{ch:hotchocolate-quickly}'s lesson about the compiler beating the
documentation, arriving again in a place I did not expect it.

\subsection*{A type nothing returns is not in the schema}

\index{HotChocolate!unreachable types}%
This one cost me twenty minutes and it will cost you the same, so here it is
before chapter~\ref{ch:the-first-cut-extracting-catalog} needs it.

The first version of the Reviews subgraph compiled, started, answered
\code{Query.reviews} correctly, and published this:

\begin{graphqlsdl}
schema
  @link(
    url: "https://specs.apollo.dev/federation/v2.6"
    import: ["@tag", "FieldSet"]
  ) {
  query: Query
}

type Query {
  reviews: [Review!]!
  _service: _Service!
}
\end{graphqlsdl}

No \code{Product}, and no \code{\_entities} either. \code{@key} is missing from
the import list, and so is the \code{\_Entity} union. The C\# was fine: the
class was there, it had its \code{[Key("id")]}, it had its reference resolver.
What it did not have was any root field returning it, and HotChocolate builds
the types it can reach from the roots. An entity that only ever arrives through
\code{\_entities} is reachable from nowhere.

One line fixes it:

\begin{minted}{csharp}
builder.AddGraphQL()
    .AddApolloFederation()
    .AddWireReviews()
    .AddType<Product>();
\end{minted}

What makes this worth a section is how quietly it fails. Nothing throws. The
subgraph is healthy and answers its own queries. The schema it publishes is
valid, and it \emph{composes}: you get a supergraph in which \code{Product} has
a title and a price and no reviews at all, and the first sign of trouble is a
client asking for a field the router says does not exist. If a subgraph you
extracted has mysteriously stopped contributing, this is the first thing to
check, and \code{\_service} is where to check it.

\index{subgraph!contract|)}%
